\section{Results}
\subsection{A broadening mean layer with a high-rank stochastic background}
The mean density layer broadens continuously over the low-Knudsen interval (\cref{fig:meanrms}). The half-jump centre changes only modestly between $\KnD=0.01$ and $0.025$, but $\deltaten/R$ nearly doubles. At $\KnD=0.05$ the transition spans a substantial fraction of the radius. In contrast, the raw density RMS remains distributed over the complete compression region and is not concentrated on a single deterministic shape.

\begin{figure}[t]
\centering
\includegraphics[width=\textwidth]{fig03_mean_rms_contours.pdf}
\caption{Mean compression-layer geometry and raw fluctuation level. Top row: $\mean\rho/\rho_\infty$ with the mean half-jump centre overlaid. Bottom row: $\rms\rho/\rho_\infty$. The same colour range is used across cases in each row. The mean layer broadens strongly while the stochastic field remains spatially distributed.}
\label{fig:meanrms}
\end{figure}

This high dimensionality is confirmed by the corrected physical-domain POD. \Cref{fig:highrank} compares the mean geometry with field compactness across all nine Knudsen numbers. For the combined state, $E_1$ decreases from 4.32\% at $\KnD=0.01$ to 1.10\% at $\KnD=1$, and the number of modes required for 90\% variance increases from 146 to 172. Density, pressure, Mach number and both temperatures follow the same qualitative conclusion. The intermediate points do not reveal modal condensation near $\KnD=0.1$; they form a smooth high-rank sequence.

\begin{figure}[t]
\centering
\includegraphics[width=\textwidth]{fig04_geometry_highrank.pdf}
\caption{Mean geometry and corrected field dimensionality over the nine-case Knudsen sweep. (a) Sector-median half-jump location and 10--90 width. (b) Corrected leading POD energy. (c) Number of modes required for 90\% variance. Exactly 200 snapshots are used at every $\KnD$ for the POD comparison.}
\label{fig:highrank}
\end{figure}

\begin{table}
\centering
\caption{Corrected common-200 POD metrics for the combined density--Mach--temperature--pressure state. $C_{10}$ is cumulative energy through ten modes.}
\label{tab:pod}
\small
\begin{tabular}{cccc}
\toprule
$\KnD$ & $E_1$ (\%) & $C_{10}$ (\%) & $N_{90}$ \\
\midrule
0.010 & 4.32 & 21.91 & 146 \\
0.025 & 2.34 & 14.87 & 157 \\
0.050 & 1.63 & 12.12 & 162 \\
0.075 & 1.81 & 11.84 & 164 \\
0.100 & 1.45 & 11.08 & 165 \\
0.150 & 1.47 & 10.46 & 167 \\
0.250 & 1.20 & 9.26 & 169 \\
0.500 & 1.24 & 8.99 & 170 \\
1.000 & 1.10 & 8.45 & 172 \\
\bottomrule
\end{tabular}
\end{table}

The collective coordinate identified below is therefore not the leading field POD mode. It is a low-amplitude covariance direction embedded in a broadband background. This distinction prevents the common but incorrect inference that a coherent marker necessarily implies a low-rank complete flow field.

\subsection{Coarse graining reveals a persistent low-Knudsen component}
\Cref{fig:coarse} shows how marker variance changes when density fields are averaged before extraction. The pointwise centre variance decreases strongly with $m$ in all cases, confirming that raw marker motion is dominated by finite-particle fluctuations. The angular-mean variance behaves differently. At $\KnD=0.01$ and $0.025$, the fitted physical contribution attenuates on a finite time scale and remains distinguishable from the sampling component. At $\KnD=0.05$, a nominal two-component fit can be drawn, but it fails the out-of-sample and PSD-correction tests below.

\begin{figure}[t]
\centering
\includegraphics[width=\textwidth]{fig05_temporal_coarse_graining.pdf}
\caption{Temporal coarse graining of the density half-jump centre. Top row: mean pointwise angular variance; bottom row: variance of the angular-mean displacement. Symbols are measurements after averaging $m$ consecutive fields before marker extraction. Solid curves are the fitted physical-plus-sampling model; dashed and dotted curves show the inferred components. The $\KnD=0.05$ decomposition is displayed as a diagnostic but is rejected by out-of-sample validation.}
\label{fig:coarse}
\end{figure}

The inference metrics are summarized in \cref{tab:inference}. At $\KnD=0.01$, the two-component model gives $\Delta AIC_c=76.4$ with a bootstrap lower bound of 33.3, and reduces the leave-one-group-size-out error from 0.784 to 0.432. At $\KnD=0.025$, the 600-snapshot record gives $\Delta AIC_c=59.1$ and reduces the error from 0.437 to 0.254. The physical angular modes have correlations 0.868 and 0.876 with uniform displacement, positive far-angle covariance, and time scales 0.653 and 0.724. Their direct shape correlation is 0.972. The same structure is therefore recovered from two independent rarefaction levels and records of different length and cadence.

\begin{table}
\centering
\caption{Correlated-noise inference for the resolved cases and the $\KnD=0.05$ diagnostic. Brackets denote 95\% moving-block bootstrap intervals. The LOOCV column is the two-component/noise-only error ratio.}
\label{tab:inference}
\small
\begin{tabular}{cccccc}
\toprule
$\KnD$ & $N_s$ & $\tau_p^*$ & $\Delta AIC_c$ & LOOCV ratio & $r_{unif}$ / $r_{far}$ \\
\midrule
0.010 & 200 & 0.653 [0.425,0.786] & 76.4 [33.3,73.2] & 0.55 & 0.868 [0.704,0.924] / 0.184 [0.065,0.331] \\
0.025 & 600 & 0.724 [0.239,1.019] & 59.1 [29.0,64.4] & 0.58 & 0.876 [0.781,0.891] / 0.267 [0.120,0.287] \\
0.050 & 419 & 0.280 [0.222,0.611] & 18.0 [1.9,37.5] & 3.74 & 0.767 [0.241,0.795] / 0.072 [-0.002,0.089] \\
\bottomrule
\end{tabular}
\end{table}

\subsection{Angular covariance and collective displacement shape}
The physical covariance matrices in \cref{fig:covariance} provide the clearest visual evidence. At $\KnD=0.01$ and $0.025$, broad positive off-diagonal regions connect rays separated by tens of degrees. Their leading modes are same-signed over the complete sector and displace the mean half-jump front with a non-uniform angular envelope. The mode should therefore be described as a collective displacement, not as rigid translation. At $\KnD=0.05$, the matrix is patchy, the PSD correction is 54\%, and the two-component LOOCV error is 3.7 times the noise-only error. That panel is retained only to show the failure of identifiability.

\begin{figure}[t]
\centering
\includegraphics[width=\textwidth]{fig06_covariance_collective_mode.pdf}
\caption{Noise-separated angular covariance and inferred displacement shape. Top row: correlation matrices of $C_p$. Bottom row: mean density half-jump front, normalized mode amplitude and $\pm1\sigma$ leading-mode displacements. The $\KnD=0.05$ result is diagnostic and is not accepted as a physical covariance mode.}
\label{fig:covariance}
\end{figure}

The angular-time maps in \cref{fig:spacetime} show why a direct visual interpretation of raw markers is unreliable. All three cases contain localized streaks and speckle. Projection onto the inferred mode exposes a sector-wide component in the two resolved cases, but a visually structured projection can also be constructed in the rejected case. The quantitative covariance gates, rather than appearance alone, determine whether the projection is physical.

\begin{figure}[t]
\centering
\includegraphics[width=\textwidth]{fig07_marker_spacetime.pdf}
\caption{Angular-time displacement maps. Top row: raw centred half-jump marker. Bottom row: projection onto the inferred leading angular mode. Each case is plotted with its own symmetric colour scale to expose structure; the panels are not direct amplitude comparisons. Quantitative model selection is required because a visually organized projection can exist even when the covariance decomposition fails.}
\label{fig:spacetime}
\end{figure}

\subsection{The same coordinate is present in the complete macroscopic fields}
The full-field matched filter tests whether the ray marker corresponds to translation of the entire compression layer. \Cref{fig:multimoment} shows amplitude correlations and representative time series. At $\KnD=0.01$, density and pressure amplitudes correlate with the marker by 0.946 and 0.901. Mach number and translational temperature give 0.755 and 0.707. A principal component formed from the four field amplitudes contains 87.3\% of their variance and correlates 0.886 with the marker. At $\KnD=0.025$, density and pressure remain strong at 0.790 and 0.727; the four-moment consensus contains 70.6\% and correlates 0.647 with the marker. At $\KnD=0.05$, the density correlation falls to 0.183 and cross-variable consensus weakens.

\begin{figure}[t]
\centering
\includegraphics[width=\textwidth]{fig08_multimoment_validation.pdf}
\caption{Full-field validation. Top row: correlation matrices of the marker and translation-template amplitudes extracted from density, Mach number, translational temperature and pressure. Bottom row: standardized marker, density and pressure amplitudes. The resolved cases show a shared multi-moment time dependence; the $\KnD=0.05$ case is substantially weaker.}
\label{fig:multimoment}
\end{figure}

\begin{table}
\centering
\caption{Full-field displacement-template metrics. Each entry is marker correlation followed by the median fraction of raw instantaneous field variance aligned with the template. Correlation intervals are given in the text.}
\label{tab:template}
\small
\begin{tabular}{ccccc}
\toprule
$\KnD$ & $\rho$ & $M$ & $T_{tr}$ & $p$ \\
\midrule
0.010 & 0.946 / 5.29\% & 0.755 / 2.30\% & 0.707 / 1.24\% & 0.901 / 4.03\% \\
0.025 & 0.790 / 2.37\% & 0.395 / 1.73\% & 0.273 / 0.97\% & 0.727 / 1.60\% \\
0.050 & 0.183 / 0.67\% & 0.325 / 1.06\% & 0.164 / 0.50\% & 0.510 / 0.93\% \\
\bottomrule
\end{tabular}
\end{table}

The translation coordinate is weak in an energetic sense. Even at $\KnD=0.01$, the median projected fraction is 5.3\% for density and 4.0\% for pressure; it is smaller for the thermal and Mach-number fields. This is not a defect of the result. It demonstrates that the coherent motion coexists with high-dimensional kinetic fluctuations and explains why conventional field POD does not isolate it as a dominant mode.

The full-field agreement is not a tautological consequence of using density to define the marker. First, pressure, Mach number and translational temperature are projected independently using their own mean gradients; their temporal amplitudes are not used in the marker extraction. Second, circularly shifting the marker relative to each field gives null-test probabilities of 0.0054 for all four moments at $\KnD=0.01$ and for density, Mach number and pressure at $\KnD=0.025$; the density probability at $\KnD=0.05$ rises to 0.0162 and its bootstrap correlation interval approaches zero. Third, the spatial template is shifted by $\Delta\xi=\{-1,-0.5,0.5,1\}$. At $\KnD=0.01$, the unshifted density template captures 5.29\% of variance, compared with 1.25--2.39\% for upstream shifts and 0.29--1.99\% for downstream shifts. At $\KnD=0.025$, the unshifted density value is 2.37\%, approximately twice the $\Delta\xi=-0.5$ value. Pressure exhibits the same zero-shift maximum in both resolved cases. At $\KnD=0.05$, by contrast, the density projection is larger for $\Delta\xi=0.5$ and 1 than for the physical zero-shift template. The resolved low-Knudsen signal is therefore localized on the actual mean compression gradient, whereas the diagnostic $0.05$ response is not.

Conditional composites improve the signal-to-noise ratio without changing the template. \Cref{fig:reconstruction} subtracts the mean of the most negative marker events from the mean of the most positive events. At both resolved Knudsen numbers, the actual density composite has a same-signed front-normal dipole over a broad angular sector. The translation reconstruction captures the global front structure, with weighted composite $R^2=0.922$ and 0.737. The residual retains fine angular and downstream variation but no comparable sector-wide density ridge.

\begin{figure}[t]
\centering
\includegraphics[width=\textwidth]{fig09_density_template_reconstruction.pdf}
\caption{Conditional density displacement reconstructed from the full field. The left column is half the difference between composites conditioned on the upper and lower 15\% of marker amplitude, the middle column is the fitted translation template, and the right column is the residual. The two rows correspond to $\KnD=0.01$ and $0.025$.}
\label{fig:reconstruction}
\end{figure}

\subsection{Consistency under early mean-layer broadening}
The principal new separation between mean and dynamic behavior is shown in \cref{fig:dynamic}. Between $\KnD=0.01$ and $0.025$, the sector-median 10--90 density width increases from $0.123R$ to $0.224R$, an $82\%$ broadening, whereas the mean half-jump standoff changes by only $1.3\%$. The layer therefore undergoes substantial internal spreading before its centre moves appreciably. In spite of this static change, the inferred angular displacement modes remain almost collinear: their absolute normalized inner product is $0.972$ (the demeaned Pearson shape correlation is $0.885$). Both modes remain same-signed over the complete $120^\circ$--$179^\circ$ sector, although the angular envelope is not exactly uniform. Early rarefaction thus changes the thickness of the compression layer much more strongly than it changes the geometry of the resolved collective motion.

The time-scale comparison in \cref{fig:dynamic}(b) strengthens this interpretation. The inferred memories are $\tau_p^*=0.653$ and $0.724$, while the corresponding layer-passage times are only $t_\delta^*=0.0615$ and $0.112$. Hence the covariance persists for $10.6$ and $6.5$ layer passages. It also persists for $2.83$ and $3.09$ standoff passages. A marker perturbation generated only by local advective crossing or by independent cell noise would not naturally retain a common angular sign for so many geometric transit times. The slow memory supports an interpretation in which the bow layer samples a global, body-coupled displacement direction that is continuously forced by the stochastic kinetic background.

\begin{table}
\centering
\caption{Dynamic consistency and multi-moment synchronization. The shape similarity is the absolute normalized inner product with the $\KnD=0.01$ mode. $S_1$ is defined in \cref{eq:synch}.}
\label{tab:dynamic}
\small
\begin{tabular}{ccccccc}
\toprule
$\KnD$ & $\deltaten/R$ & $\tau_p^*$ & $\tau_p^*/t_\delta^*$ & shape similarity & $S_1$ & median $r_{qq'}$\\
\midrule
0.010 & 0.123 & 0.653 & 10.61 & 1.000 & 0.873 & 0.807\\
0.025 & 0.224 & 0.724 & 6.46 & 0.972 & 0.706 & 0.555\\
0.050 & 0.359 & -- & -- & -- & 0.529 & 0.369\\
\bottomrule
\end{tabular}
\end{table}

\begin{figure}[t]
\centering
\includegraphics[width=\textwidth]{fig11_dynamic_invariance_and_decoherence.pdf}
\caption{Static broadening and dynamic organization. (a) Angular eigenfunctions at $\KnD=0.01$ and $0.025$, normalized by their sector means. Their absolute normalized inner product is $0.972$. (b) Collective memory compared with the mean standoff- and 10--90-layer passage times. (c) Coherent displacement gain $G_q$ of density, pressure, Mach number and translational temperature. (d) Multi-moment synchronization measures. The mean layer broadens strongly from $\KnD=0.01$ to $0.025$ while the angular displacement geometry and slow memory persist, but the synchronous participation of the non-density moments weakens. The $\KnD=0.05$ values are shown as diagnostic matched-filter outputs rather than as evidence of a resolved collective covariance.}
\label{fig:dynamic}
\end{figure}

\subsection{Moment-selective dynamic decoherence}
Persistence of the marker shape does not imply that every macroscopic moment continues to translate with equal amplitude. Equation~\eqref{eq:gain} separates this issue. At $\KnD=0.01$, the coherent gains are $G_\rho=0.828$, $G_p=0.781$, $G_M=0.702$ and $G_{T_{tr}}=0.690$. The four moments therefore participate in a nearly common displacement, with density and pressure slightly closer to the marker than Mach number and translational temperature. The first standardized multi-moment component contains $87.3\%$ of amplitude variance, its correlation with the marker is $0.886$, and the median pairwise moment correlation is $0.807$.

At $\KnD=0.025$, the displacement coordinate remains resolved but the participation hierarchy becomes much stronger. The density and pressure gains decrease to $0.631$ and $0.575$, while the Mach-number and translational-temperature gains fall to $0.367$ and $0.312$. The first multi-moment component now contains $70.6\%$ of the variance and the median cross-moment correlation decreases to $0.555$. Thus the front geometry remains collective while the thermodynamic response becomes less synchronous. This behavior is distinct from the static variable-dependent scale hierarchy documented in the companion paper: here the separation is measured in the temporal amplitudes of a common displacement template, not in parameter-dependent mean profiles.

The phase portraits in \cref{fig:phase} make this distinction visible without energy-ranking the fields. Density remains closely slaved to the marker at $\KnD=0.01$ and $0.025$, whereas the translational-temperature cloud broadens and rotates toward weak correlation at $0.025$. At $\KnD=0.05$, neither phase portrait provides a robust common-coordinate signature under the full inference criteria. The physically conservative conclusion is therefore \emph{moment-selective weakening of synchrony}: early rarefaction preserves a slow compression-coordinate geometry, but velocity--sound-speed and thermal adjustments increasingly acquire independent fluctuations and relaxation content.

\begin{figure}[t]
\centering
\includegraphics[width=\textwidth]{fig12_multimoment_phase_portraits.pdf}
\caption{Standardized phase portraits of marker displacement $a_m$ against density and translational-temperature equivalent displacement amplitudes. The density relation remains narrow at $\KnD=0.01$ and $0.025$, while the thermal relation loses coherence more rapidly. The $\KnD=0.05$ panels are diagnostic matched-filter outputs; the covariance mode is not classified as resolved there.}
\label{fig:phase}
\end{figure}

\subsection{What can be concluded beyond $\KnD=0.025$}
A non-detection is meaningful only when the analysis could have detected a mode of relevant size. Sliding-window and injection tests are summarized in \cref{fig:limits}. The $\KnD=0.025$ mode passes the strict window criterion in six of nine independent overlapping windows. At $\KnD=0.1$, only two early windows pass, whereas the complete record is ambiguous. No windows pass at $\KnD=0.05$, $0.075$ or $0.15$ under the combined criterion.

The injection test uses the $\KnD=0.01$ shape and time scale with the measured noise covariance of each case. At $\KnD=0.025$, detection power at the reference amplitude is 0.89 and the interpolated 90\% threshold is $U_{90}/A_{ref}=1.03$. For $\KnD=0.05$ and $0.075$, the thresholds rise to 5.0 and 5.33. At $\KnD=0.1$ and $0.15$, 90\% power is not reached even at six times the reference amplitude. Consequently, the data beyond $\KnD=0.025$ cannot exclude persistence of a low-Knudsen-sized mode. The correct result is resolved emergence at two near-continuum states, not demonstrated disappearance with rarefaction.

\begin{figure}[t]
\centering
\includegraphics[width=\textwidth]{fig10_identifiability_limits.pdf}
\caption{Persistence and detectability limits. (a) Sliding-window $\Delta AIC_c$; black stars denote windows passing the combined model-selection, LOOCV, shape, far-angle and PSD criteria. (b) Detection power after injecting the $\KnD=0.01$ collective mode into the measured noise of each case. (c) Amplitude required for 90\% detection relative to the reference low-Knudsen amplitude. Values above unity mean that the record cannot exclude a mode as small as the resolved reference.}
\label{fig:limits}
\end{figure}
